% Generated draft: requires mathematical and source review.
\documentclass[11pt]{article}
\usepackage[margin=25mm]{geometry}
\usepackage{fontspec}
\setmainfont{Latin Modern Roman}
\usepackage{amsmath,amssymb,mathtools,unicode-math}
\setmathfont{Latin Modern Math}
\usepackage{enumitem,needspace,xurl,hyperref}
\setlength{\parindent}{0pt}
\setlength{\parskip}{5pt}
\setlength{\emergencystretch}{3em}
\newcommand{\N}{\mathbb{N}}
\newcommand{\Z}{\mathbb{Z}}
\newcommand{\Q}{\mathbb{Q}}
\newcommand{\R}{\mathbb{R}}
\newcommand{\C}{\mathbb{C}}
\newcommand{\F}{\mathbb{F}}
\newcommand{\A}{\mathbb{A}}
\newcommand{\PP}{\mathbb{P}}
\newcommand{\E}{\mathbb{E}}
\newcommand{\eps}{\varepsilon}
\newcommand{\dd}{\,\mathrm{d}}
\begin{document}
\section*{UnsolvedMath ID: 10400186}
\textbf{Source problem code:} catalogue entry
\subsection{Definitions and mathematical statement}
Let \(\Lambda\) denote Vogel's algebra with the coefficient, diagram, relation, grading and multiplication conventions of the source's surrounding discussion. Those defining conventions are absent from this excerpt. The task is to describe \(\Lambda\), for example by giving a complete set of generators and relations.

More precisely, after its ground ring \(k\) and algebra category are fixed, a presentation consists of specified elements \(\lambda_s\in\Lambda\), indexed by a set \(S\), and explicit relations \(\mathcal R\) in the corresponding free algebra \(F_k(S)\), such that the homomorphism
\[\pi:F_k(S)\longrightarrow\Lambda,\qquad s\longmapsto\lambda_s,\]
is surjective and
\[\ker\pi=(\mathcal R),\qquad \Lambda\cong F_k(S)/(\mathcal R).\]
Here \((\mathcal R)\) is the appropriate generated ideal; the free algebra must be chosen to match whether the supplied algebra is commutative, unital and graded. Surjectivity expresses completeness of the generators, and the kernel equality expresses completeness of the relations. Any claimed grading compatibility must use the source's grading. The problem does not demand that the presentation be finite, or claim that finite presentation is possible. The phrase ``say'' makes generators and relations one suggested form of a full description, rather than the only allowed kind of answer. No diagram definition of \(\Lambda\) or list of purported Vogel generators is reconstructed without the omitted conventions.
\subsection{Short English statement}
Describe Vogel's algebra, for example through a presentation that generates every element and accounts for every relation. The source does not require finitely many generators or relations and does not reproduce the algebra's definition.
\subsection{Sources}
\begin{itemize}
\item Local UnsolvedMath registry, permanent ID \(10400186\), exact source code \texttt{catalogue entry}.
\item Ohtsuki, \emph{Problems on invariants of knots and 3-manifolds} (2002), Problem 10.11, PDF page 149; the named algebra is attributed to Vogel: \url{https://msp.org/gtm/2002/04/gtm-2002-04-024s.pdf}. Registry access 2026-07-29; extraction pdf-font-structure.
\end{itemize}
\end{document}
